% REVIEW DRAFT: not an input of the current published manuscript.
% This file is not proof certification; see geometry-audit-2026-11-21.json.

\section{The actual theta comparison and its two rank recursions}
\label{sec:sr-theta-comparison-endgame}

We now combine the proved local calculations and the global
constant-term formulas. In particular we prove a comparison of the
actual periods, with their fixed models and Haar measures.
All spaces and evaluation maps below are those constructed in
Proposition~\ref{prop:sr-theta-coherent-tower}.

Write \(\nu=\alpha/\beta\), \(T=S_0\cup Q\), \(d=a+a'\), and
\(\delta=(a-a')/2\). If an adapted first plane and target plane
have coefficient ideals \(\mathfrak b,\mathfrak b'\), define
\[
 \begin{split}
 A_\nu&=\nu(\mathfrak b)\nu(\mathfrak b')
       \prod_{v\in T}\epsilon_v(\nu,1-\delta),\\
 B_\nu&=\{\nu(\mathfrak b)\nu(\mathfrak b')\}^{-1}
       \prod_{v\in T}\epsilon_v(\nu,-\delta).
 \end{split}                                      \tag{E1}
 \label{eq:sr-theta-two-step-factors}
\]
The ideal values and epsilon factors have exactly the conventions
of the finite standard-vector calculation. A place of \(T\)
has the specified split embedding; the epsilon factor in (E1)
is for its chosen factor of \(E_v=F_v\times F_v\).
There is one factor for each \(F\)-place of \(T\), as in that
calculation. In particular neither a second split factor nor
an additional sign is inserted.

\begin{srlemma}[Comparison of the actual lower quotient maps]
\label{lem:sr-theta-actual-lower-comparison}
Suppose both lower periods for a rank reduction are nonzero and
proportional, with the actual matching quotients.
For an \(A\)-reduction with \(q=0\) there is a scalar
\(k_A\ne0\) such that every upper polynomial matching pair
\((\Phi,\Phi')\) satisfies the lower polynomial matching relation
\[
     C_A\Phi
       \ \sim_1\
       k_A A_\nu L^{S_0}_E(\nu,\delta)\,e'_0\Phi'.
 \tag{E2}\label{eq:sr-theta-actual-LocA}
\]
For a \(B\)-reduction there is a scalar \(k_B\ne0\) such that
\[
   k_B B_\nu L^{S_0}_E(\nu,\delta+1)\,e_0\Phi
                  \ \sim_1\ C'_B\Phi'.           \tag{E3}
 \label{eq:sr-theta-actual-LocB}
\]
Here \(C_A,C'_B\) are the actual global cone operators,
\(e_0,e'_0\) are the actual global zero evaluations, and
\(\sim_1\) means equality in the lower common tensor quotient.
In (E3), if \(\Phi'\) is in the actual upper full-rank
Schwartz subspace at \(v_1\), then \(C'_B\Phi'\) belongs
to the corresponding lower subspace.

The scalars \(k_A,k_B\) depend on the fixed ground, geometric,
measure, upper-real-quotient and lower-real-quotient data.
They are chosen before the finite packet \(T,\alpha_f,\beta_f\)
and before all finite matching quotient presentations.
\end{srlemma}

\begin{proof}
Proposition~\ref{prop:sr-theta-constant-term-descent} gives
the actual descended local maps on the upper common quotient,
on both sides of the comparison. Its nonvanishing hypotheses
are precisely the lower nonvanishing assumed here.
We compare these descended maps place by place.

At a finite place outside \(T\), the upper representation is
generated by its lattice vacuum. Iwasawa decomposition
\(G_v=P_vK_v\) implies that any \(P_v\)-intertwining map
is determined by its value on that vector.
Both descended maps send it into the one-dimensional lower
spherical line. Zero evaluation sends it to the lower vacuum,
and the cone calculation sends it to
\(u_v\) times that vacuum, where
\[
 u_v=\prod_{w\mid v}
       \frac{(\nu(\varpi_w)Nw^{-z})^{\operatorname{ord}_w\mathfrak a}}
                     {1-\nu(\varpi_w)Nw^{-z}},
 \quad
 \mathfrak a=
 \begin{cases}
  \mathfrak b\mathfrak c\,\overline{\mathfrak b'}^{-1},&A,\\
  \mathfrak b^{-1}\overline{\mathfrak b'},&B.
 \end{cases}                                      \tag{E4}
 \label{eq:sr-theta-endgame-unramified-scalar}
\]
For \(A\), \(z=\delta\); for \(B\), \(z=\delta+1\).
Every denominator and numerator here is nonzero.
These are the previously calculated complete geometric-series
integrals. This comparison applies also to split dyadic places.
At \(v_1\), where the small test space has full-rank support,
the cone augmentation identity first extends its descended
map to the full quotient. The full-rank cutoff represents
the vacuum's quotient class. Thus the same vacuum comparison
determines the map at \(v_1\); it does not require the vacuum
itself to have full-rank support.

At \(v\in T\), the ramified branching calculation proves that
the two nonzero descended \(P_v\)-maps are proportional.
Its standard matched pair has nonzero lower zero evaluation
and satisfies, exactly,
\[
 \begin{array}{ll}
  p_{1,v}(\Lambda_{A,v}\phi)
    =\epsilon_v(\nu,1-\delta)\,
                         p'_{1,v}(e'_{0,v}\phi'),& A,\\
  p'_{1,v}(\Lambda_{B,v}\phi')
    =\epsilon_v(\nu,-\delta)\,
                         p_{1,v}(e_{0,v}\phi),& B.
 \end{array}                                      \tag{E5}
 \label{eq:sr-theta-endgame-ramified-scalar}
\]
The one-dimensional branching space therefore extends (E5)
from this standard pair to every upper matched pair.
The factors are nonzero: the ramified primitive Gauss sum is
nonzero and the remaining conductor powers are units.
This argument fixes the relative map, independently of a
common rescaling or another presentation of the finite
matching quotient.

At infinity the two descended maps are in the single branching
line proved from the actual bottom compact type.
Zero evaluation is nonzero because it is onto the lower
polynomial space. The cone map is nonzero by the actual
cone bump construction together with lower injectivity.
Consequently their proportionality scalar \(c_\infty\)
is nonzero. It is one scalar for all the finite packets,
for the following reason. The infinite characters
\(\alpha_\infty,\beta_\infty\) are determined by the fixed
weights of the ground datum. The infinite Weil action and
the infinite cone integral therefore remain the same.
The upper and lower polynomial quotient maps have also
been fixed before the packet. Thus the two concrete maps
whose ratio is \(c_\infty\) are fixed maps. Existence of a
packet was used to prove their descent, but cannot change
that ratio. If there is no allowable packet, choose the
nonzero scalar \(1\) for the vacuous assertion.

We spell out why these local statements imply (E2)--(E3)
for a tensor matching relation, rather than merely for a
pure pair. Take a finite set \(U\supset T\cup\{v_1\}\)
outside which all the finite test functions are vacua.
Proposition~\ref{prop:sr-cone-euler-factorization} expresses
the global cone as
\[
 \kappa\,\nu(\mathfrak a)\,L^{S_0}_E(\nu,z)
       \left(\prod_{v\in U\setminus T}u_v\right)^{-1}
       L_\infty\otimes\bigotimes_{v\in U}\Lambda_v,
 \tag{E6}\label{eq:sr-theta-endgame-tensor-cone}
\]
with the lower vacua outside \(U\).
Its nonzero \(\kappa\) includes the fixed Haar and ideal-norm
factors. The finite-place map comparison contributes
\(\prod_{v\in T}\epsilon_v\prod_{v\in U\setminus T}u_v\);
the second product cancels the inverse in (E6).
The infinite comparison contributes \(c_\infty\).
Finally the computed coefficient ideals give
\[
 \nu(\mathfrak a)=
 \begin{cases}
   \nu(\mathfrak b)\nu(\mathfrak b'),&A,\\
   \{\nu(\mathfrak b)\nu(\mathfrak b')\}^{-1},&B,
 \end{cases}
\]
since \(\nu(\mathfrak c)=1\) and
\(\nu(\bar{\mathfrak b'})=\nu(\mathfrak b')^{-1}\).
This proves the stated factors, with \(k=\kappa c_\infty\).
For a finite sum, apply these linear maps to the equality
of its tensors in the upper common quotient.
No termwise pure matching is assumed. Enlarging \(U\)
only inserts another \(u_v/u_v\), so the answer is independent
of the chosen finite presentation.

For \(B\), the explicit rank calculation of the local cone
operator at \(v_1\) preserves the lower full-rank cutoff.
At every other place its output is Schwartz--Bruhat and at
infinity it is polynomial Gaussian. The restricted product
formula (E6) and linearity prove the claimed global cutoff
membership. This completes both the actual matching
relations and the asserted order of choices.
\end{proof}

\begin{srtheorem}[The actual common theta-period image]
\label{thm:sr-theta-actual-proportionality}
For every specified triple in the actual rank and weight range,
and for every actual real and finite matching quotient system,
there is a unique nonzero \(c\in\mathbb C\) such that
\[
            I_V(\Phi)(h)=c\,I_{V'}(\Phi')(h)
 \tag{E7}\label{eq:sr-theta-actual-proportionality}
\]
for every matching pair and every \(h\in U(W)(\mathbb A_F)\).
Both polynomial period images are nonzero; the small image
already is nonzero on the specified full-rank Schwartz subspace.
\end{srtheorem}

\begin{proof}
We induct on \(Y=p+q\), simultaneously for all actual ranks,
lattices, models, and matching quotient systems satisfying
the hypotheses. Proportionality at one node is therefore
available for every lower node constructed above.

For \(Y=0\), both first spaces and the target are definite.
The compact global comparison proves (E7), with nonzero scalar,
using the explicitly normalized anisotropic Siegel--Weil
identity and its nonzero standard sections.
The small full-rank quotient is onto the common representation,
so Factor supplies a full-rank representative with nonzero
small period. This gives the required base nonvanishing.

Suppose \(Y\ge1\). Factor realizes the two polynomial periods,
the small one on its full-rank subspace, as maps from the
same irreducible restricted tensor quotient \(\Pi\).
We first compare on these polynomial vectors.

If \(q\ge1\), take a \(B\)-reduction and invoke induction
on its lower pair. Equation (E3) and the lower comparison
constant \(c_1\), combined with the exact \(\beta\) constant
terms of the two upper periods, give a nonzero number \(c\)
such that the difference
\[
                  F_\xi=I_V(\xi)-cI_{V'}(\xi)
\]
has zero \(\beta\)-component for every \(\xi\in\Pi\).
All factors used in solving for \(c\) are nonzero:
\(k_B\), the primitive epsilon factors, the ideal values,
\(c_1\), and \(L_E^{S_0}(\nu,\delta+1)\).
The latter is an absolutely convergent unitary Hecke Euler
product in \(\operatorname{Re}s>1\).
Proposition~\ref{prop:sr-theta-two-step-zero} makes its
\(\alpha\)-component zero as well. If \(p=0\), its boundary
case gives this directly; if \(p\ge1\), its rank-one
contraction uses the \(A\)-lower node and the further
\(B\)-lower node, with signature sums \(Y-1\) and \(Y-2\).
Thus it invokes only the induction hypotheses.

If \(q=0\), then \(p\ge1\). Take an \(A\)-reduction.
Equation (E2), its lower comparison, and the exact \(\alpha\)
constant terms give a nonzero \(c\) for which
\(F_\xi\) has zero \(\alpha\)-component.
The Euler product here is \(L_E^{S_0}(\nu,\delta)\), again
absolutely convergent and nonzero.
The boundary vanishing already proved makes its
\(\beta\)-component zero.

In both cases the difference has identically zero line
constant term. Proposition~\ref{prop:sr-theta-line-implies-cusp}
makes it cuspidal. It lies in the span of the two actual
period images, so
Corollary~\ref{cor:sr-theta-cuspidal-difference-zero} makes
it zero. For isotropic first spaces that corollary uses
the full-rank rapid-decay backlift. If one first space
is definite, its proof uses the anisotropic doubled
kernel and the nonzero period in the appropriate
opposite rank reduction. That lower node has signature
sum \(Y-1\), already covered by induction.
This proves (E7) on the common polynomial quotient.

We next check nonvanishing at this node; it was not assumed.
In the \(B\)-case, the large-side zero evaluation has a
nonzero lower quotient image: choose a lower polynomial
vector and lift at infinity, choose the standard ramified
vectors at \(T\), and the vacua elsewhere.
The resulting descended map is nonzero and is equivariant
for the lower group and its Lie/compact actions. Its image
is therefore the whole irreducible lower quotient.
As the lower period map is injective and nonzero,
the upper large \(\beta\)-component is nonzero on some vector.
Thus \(I_V\ne0\), and (E7), with \(c\ne0\), gives
\(I_{V'}\ne0\) on the small full-rank polynomial subspace.

In the \(A\)-case the small first space is definite.
The small-side zero evaluation onto its lower polynomial
space supplies a vector with nonzero lower period.
At \(v_1\) it can be lifted with full upper row rank:
in the split coordinate formula choose a compactly supported
bump in the integrated columns whose integral is one;
the already full-rank lower row block stays full rank after
these columns are adjoined. Hence it has a representative
in the actual upper full-rank subspace.
The upper small \(\alpha\)-component is nonzero.
Equation (E7) then makes the upper large period nonzero.
This completes the simultaneous induction and its
nonvanishing assertion without circularity.

It remains to treat the precise matching notion for general
Schwartz vectors. Its definition supplies polynomial matching
approximations with fixed finitely many finite test functions,
with the small approximations in the specified full-rank
subspace. The integrated Schwartz-seminorm bounds proved for
the large period and for the small full-rank period imply
continuity on each such finite test space at each fixed \(h\).
They were obtained by retaining the Schwartz seminorm in the
coefficient-square and full-rank minor majorants; thus this
passage uses a genuine seminorm estimate, not constants
depending arbitrarily on each vector.
Both sides of (E7) consequently converge to the periods of
the stated matching pair. This proves (E7) for that pair
and every \(h\).

Finally take a nonzero small polynomial period with full-rank
support, and lift its common quotient tensor to a large
polynomial test vector. They form a matching pair with
constant polynomial approximation. Evaluating where the
small period is nonzero determines \(c\) uniquely.
The established nonvanishing of the large common-quotient
image excludes \(c=0\).
\end{proof}

\begin{srproposition}[The two recursion constants, with packet quantifiers]
\label{prop:sr-theta-actual-StepAlphaBeta}
For each fixed adapted rank reduction and fixed upper and
lower real matching maps there are constants \(C_A,C_B\ne0\),
chosen before the finite packet and all finite matching maps,
such that the unique actual comparison constants satisfy
\[
 \begin{split}
 c&= C_A\,
       \nu(\mathfrak b)\nu(\mathfrak b')
       \prod_{v\in T}\epsilon_v(\nu,1-\delta)\,
       L_E^{S_0}(\nu,\delta)\,c_1
                     &&(A,\ q=0),\\
 c\,L_E^{S_0}(\nu,\delta+1)
    &=\left\{C_B\,
       [\nu(\mathfrak b)\nu(\mathfrak b')]^{-1}
       \prod_{v\in T}\epsilon_v(\nu,-\delta)
      \right\}^{-1}c_1
                     &&(B).
 \end{split}                                      \tag{E8}
 \label{eq:sr-theta-actual-recursions}
\]
In particular the integer-unit exponents of the two ideal
values are \(+1,+1\) for \(A\), and \(-1,-1\) for \(B\);
the additional local unit sign is \(+1\) in both recursions.
\end{srproposition}

\begin{proof}
We keep the fixed nonzero geometric constants in the
four constant-term formulas. In case \(A\), let the
large \(\alpha\)-component be \(g_A I_1(C_A^{\rm cone}\Phi)\)
and the small \(\alpha\)-component be
\(g_A'I_1'(e'_0\Phi')\), with \(g_A,g_A'\ne0\).
Equation (E2) and lower proportionality turn the first
one into
\[
   g_A k_A A_\nu L_E^{S_0}(\nu,\delta)c_1\,
                                  I_1'(e'_0\Phi').
\]
Upper proportionality says this equals
\(c\,g_A'I_1'(e'_0\Phi')\).
There is a matched pair for which this lower period is nonzero,
as in the nonvanishing proof above. Canceling it gives
the first formula in (E8), with \(C_A=g_Ak_A/g_A'\).

In case \(B\), write the large \(\beta\)-component as
\(g_B I_1(e_0\Phi)\) and the small one as
\(g_B'I_1'(C_B^{{\rm cone}\,\prime}\Phi')\).
Equation (E3) gives
\[
 k_B B_\nu L_E^{S_0}(\nu,\delta+1)
       I_1(e_0\Phi)
             =c_1 I_1'(C_B^{{\rm cone}\,\prime}\Phi').
\]
Combine this with
\(g_B I_1(e_0\Phi)=
     c\,g_B'I_1'(C_B^{{\rm cone}\,\prime}\Phi')\),
and cancel a nonzero lower period. It yields the second
formula in (E8), with \(C_B=g_B'k_B/g_B\).

The \(g\)'s are fixed by the geometric Haar disintegrations,
and the \(k\)'s have exactly the pre-packet quantifiers
proved above. Thus these constants have those quantifiers
as well. They are independent of the later finite conductors,
ramification sets and finite quotient presentations.
The conductor calculation (E5) has local sign \(+1\);
the coefficient-ideal calculation gave the displayed
ideal exponents, so these are actual values of all three
integer-unit parameters rather than unnamed existential signs.
\end{proof}

We finish by recording the precise scope of the resulting
classical geometric input. The large Schwartz period and
the small full-rank Schwartz period have the exact convergence,
rapid first-variable decay and polynomial second-variable
height growth established earlier. At a good split place,
the open-frame coinvariant calculation gives both kernels
as the first-group augmentation subspaces, gives a full-rank
representative for every small quotient vector, and lets every
linear augmentation-annihilating functional factor uniquely
through that quotient. These are algebraic consequences of
the proved full-rank quotient isomorphism.
The full-rank condition is stable under the entire Weil pair:
at the distinguished split place both first and target actions
are invertible left/right matrix changes in the specified
bar-coordinate realization; at other places the operators
preserve Schwartz--Bruhat functions and at infinity they
preserve Schwartz functions. Almost all finite vacua are fixed.
Hence the actual global full-rank Schwartz subspace is stable
under every adelic pair action.

Together with the finite and real matching construction,
Theorem~\ref{thm:sr-theta-actual-proportionality} and
(E8) therefore supply the actual \(\mathrm{ThetaIntegralsAgree}\)
and \(\mathrm{LocalMatching}\) assertions, including all their
matching, nonvanishing, convergence and recursion clauses.
This statement concerns the theta part only. The
covolume and Hodge-metric construction is a separate
geometric input and is not inferred from proportionality.
